\section{Pre-training 叙事卷土重来}

上一章说明整期的主线，本章进入第一个非共识。2024 年后，行业讨论大量转向 Post-training、RL、o 系列 reasoning model、Agent 应用和产品化，因此“Pre-training 的收益是否放缓”成为流行判断。广密在这里反向强调：Pre-training 还没有结束，并且仍然是打开新能力上限的最本质阶段。

\begin{figure}[H]
\centering
\includegraphics[width=0.92\textwidth]{pretrain-posttrain-split.png}
\caption{Pre-training 与 Post-training：预训练决定内在上限，后训练和 RL 更像能力塑形。自制概念图，依据 00:04:22--00:11:56 对谈内容整理。}
\end{figure}

\begin{knowledgebox}{读图：上限和塑形不是同一个问题}
左侧 Pre-training 负责大规模压缩世界知识、代码、推理模式和潜在能力；右侧 Post-training/RL 负责把能力对齐到任务、偏好、工具和安全边界。广密的判断是：后训练可以让模型“更会用已有能力”，但很难凭空创造底座没有的能力。
\end{knowledgebox}

\subsection{为什么说 Pre-training 仍然是非共识}

本节拆解“非共识”的来源。共识的一面是，很多人看到 GPT-4 以后模型发布节奏变慢、数据墙讨论增多、RL/reasoning benchmark 快速刷分，于是认为底座预训练红利已经接近尾声。非共识的一面是，广密认为领先者阶段性放慢不等于范式结束；如果 OpenAI 的 Pre-training 团队动荡，外界可能会把一家公司的组织状态误读成整个路线的技术极限。

\begin{importantbox}{广密强调：领先者不等于路线本身}
当某个领先公司在某条线上变慢，行业容易把它解释为“这条技术路线结束了”。但也可能只是组织能力、人才流动、资源分配和战略优先级改变。判断技术路线是否结束，必须看多家公司、多种架构和下一代模型的综合表现。
\end{importantbox}

\noindent\begin{tabularx}{\textwidth}{p{0.20\textwidth}p{0.34\textwidth}X}
\toprule
判断层次 & 容易误判的地方 & 更稳妥的读法 \\
\midrule
公司进度 & OpenAI 某阶段 Pre-training 放慢 & 可能是组织和战略问题，不等于全行业放慢。 \\
Benchmark & reasoning model 快速刷分 & 刷分很快，但不必然打开全部能力上限。 \\
数据墙 & 高质量互联网文本变少 & 合成数据、代码环境、工具反馈可能改变数据供给。 \\
应用热度 & Agent/应用更吸睛 & 应用爆发仍然依赖底座智能继续上拱。 \\
\bottomrule
\end{tabularx}

\subsection{Post-training 和 RL 能做什么，不能做什么}

本节把 Post-training/RL 的角色讲清楚。广密不是否定后训练，而是反对把后训练当作底座能力的替代品。后训练可以让模型更遵循指令、更安全、更擅长特定工具、更会在数学和代码任务上展开思考；RL 可以通过奖励信号激发模型已有的推理模式。但如果 base model 的知识、表征和行动先验不足，后训练容易像“小学生刷题”：短期分数上升，长期天花板有限。

\begin{knowledgebox}{术语消化：Pre-training、Post-training、RL}
\noindent\begin{tabularx}{\textwidth}{p{0.18\textwidth}p{0.36\textwidth}X}
\toprule
术语 & 机制 & 本期中的作用 \\
\midrule
Pre-training & 在大规模数据上训练 base model，学习通用表征和预测能力 & 决定模型内在上限，是广密重新强调的主线。 \\
Post-training & 用指令、偏好、安全和任务数据塑形模型行为 & 提高可用性，但主要激发和组织已有能力。 \\
RL & 通过奖励信号优化行为，常用于 reasoning、代码、工具任务 & 能强化探索和解题模式，但依赖底座能力。 \\
Synthetic Data & 由模型或系统生成的训练数据 & 可能缓解数据墙，尤其是高价值 CoT/工具轨迹。 \\
\bottomrule
\end{tabularx}
\end{knowledgebox}

\begin{warningbox}{常见误区：把“更会考试”误认为“更聪明”}
Reasoning model 在数学、代码和 benchmark 上进步很快，但这不等于它已经获得了所有新能力。需要区分“已有能力被更好地调用”和“底层表征真的升级”。
\end{warningbox}

\subsection{合成数据和训练框架的难点}

本节进一步解释为什么 Pre-training 还可能继续。广密提到，高价值 CoT 数据、RL 生成数据和环境中的 sampling 可能重新进入预训练阶段，缓解 data bottleneck。但这不是简单把生成文本塞回训练集：它要求训练框架把 inference、sampling、reward、filtering 和 pre-training 更紧密地结合起来。换句话说，未来的预训练可能不再是静态语料训练，而是带有探索、生成和反馈的数据工厂。

\begin{importantbox}{实践经验：数据墙不只靠“更多网页”解决}
如果未来的高价值数据来自模型探索、代码执行、工具调用和环境反馈，那么数据工程会从“清洗互联网文本”转向“运行环境、采样轨迹、评价结果、回灌训练”。这也是 Agent 和 Online Learning 会回到模型训练主线的原因。
\end{importantbox}

\subsection{本章小结}

Pre-training 叙事回归，不是说后训练不重要，而是说后训练不能替代底座上限。2025 年 Q1 的关键问题，是哪些公司仍然能持续扩大 base model 能力，哪些公司只是把已有能力包装成更好看的产品和 benchmark。

